\subsection*{一步一步建立直观}

\paragraph{1. 把底面测量与高度测量分开}

先用右图的两根向量确定底面，并用叉积法线记录有向面积。若再给一根第三边 $w$，只有它沿法线的分量贡献高度。沿底面平移顶部只会剪切整个平行六面体，不会改变体积。



\paragraph{2. 两个测量合成一个标量}

点积 $w\cdot(u\times v)$ 把第三边投到法线，同时乘上底面积。它的绝对值是体积，符号表示三根边的定向。这里的叉积因而可以理解成一种测量装置：输入 $w$，返回对应的有向体积。



\paragraph{3. 退化揭示三重积为零的原因}

把底面两根边调成共线，不论第三边是什么，体积都为零。另一种零体积来自第三边落在非退化底面内。两种几何机制都被同一个标量三重积捕获；证明时应区分底面已退化与高度为零。
